\documentclass[11pt,a4paper]{article}
\usepackage[T1]{fontenc}
\usepackage[utf8]{inputenc}
\usepackage[english]{babel}
\usepackage{amsmath,amssymb}
\usepackage[margin=2cm]{geometry}
\setlength{\parindent}{0pt}
\setlength{\emergencystretch}{3em}
\allowdisplaybreaks

\title{About \texttt{false\_exponentials\_logarithms.tex}}
\author{}
\date{}
\begin{document}
\maketitle
This note explains what the test data in \texttt{false\_exponentials\_logarithms.tex} is used for. That file is input for the engine in \texttt{main.c}; it is read by the test \texttt{test\_exponentials\_logarithms.c} in this folder.

\begin{enumerate}
\item \textbf{Purpose.} The file holds 207 false statements: textbook results with one detail changed, and the one result that the textbook prints wrong.

\item \textbf{What it exercises.} The anti-regression check against false proofs. An endpoint moved to a nearby number, included or excluded by mistake, a solution dropped, a candidate outside the domain kept, no solutions claimed for an equation that has some: the engine must never prove one of them, and today it finds a counterexample for 191. One of these statements was proved by an earlier version of the engine (exercise 402, where $\log x^{2}$ stands beside $\log x$).

\item \textbf{Variety.} Every kind of change is applied to equations, inequalities and domains. Each false statement has a point where its two sides differ, found by an independent oracle (\texttt{tools/fuzz/fuzz\_exponentials\_logarithms.py}).

\item \textbf{Expected result.} The lines \texttt{PASS: 207 of 207 false statements rejected} and \texttt{PASS: 191 of 207 statements \ldots are REFUTED (at least 191 wanted), 0 are PROVED (none allowed)}.

\item \textbf{How to run it.} From the top folder of the repository:
\begin{verbatim}
cc -std=c11 -O2 -Wall -Wextra -o /tmp/test_exponentials_logarithms \
    tests/exponentials_logarithms/test_exponentials_logarithms.c \
    && /tmp/test_exponentials_logarithms
\end{verbatim}
\end{enumerate}
\end{document}
